 % 十二地支基础圆环图
 % 用途：展示十二地支的顺时针圆环排列，标注方位和季节
 \documentclass[tikz,border=10pt]{standalone}
 \usepackage{ctex}
\usepackage{xcolor}
\pagecolor{white}
 \usetikzlibrary{arrows.meta, positioning, calc, decorations.pathmorphing, backgrounds}
 
 \begin{document}
 \begin{tikzpicture}[scale=1.2]
   % 定义地支坐标（顺时针圆环，子在上=北）
   \coordinate (zi)  at ( 90:3.5);
   \coordinate (chou) at ( 60:3.5);
   \coordinate (yin)  at ( 30:3.5);
   \coordinate (mao)  at (  0:3.5);
   \coordinate (chen) at (330:3.5);
   \coordinate (si)   at (300:3.5);
   \coordinate (wu)   at (270:3.5);
   \coordinate (wei)  at (240:3.5);
   \coordinate (shen) at (210:3.5);
   \coordinate (you)  at (180:3.5);
   \coordinate (xu)   at (150:3.5);
   \coordinate (hai)  at (120:3.5);
 
   % 外圈圆环
   \draw[thick, gray!60] (0,0) circle (3.5);
 
   % 地支节点
   \foreach \name/\angle/\text in {
     zi/90/子, chou/60/丑, yin/30/寅, mao/0/卯,
     chen/330/辰, si/300/巳, wu/270/午, wei/240/未,
     shen/210/申, you/180/酉, xu/150/戌, hai/120/亥
   } {
     \fill[black] (\angle:3.5) circle (2pt);
     \node[font=\large] at (\angle:4.2) {\text};
   }
 
   % 方位标注
   \node[font=\small\bfseries, blue!60] at ( 90:2.8) {北};
   \node[font=\small\bfseries, blue!60] at (270:2.8) {南};
   \node[font=\small\bfseries, blue!60] at (  0:2.8) {东};
   \node[font=\small\bfseries, blue!60] at (180:2.8) {西};
 
   % 四季区域（浅色背景弧）
   \begin{scope}[on background layer]
     \fill[green!8]  (30:1.8) arc(30:150:1.8) -- cycle;
     \fill[red!8]    (330:1.8) arc(330:30:1.8) -- cycle;
     \fill[orange!8] (210:1.8) arc(210:330:1.8) -- cycle;
     \fill[blue!8]   (150:1.8) arc(150:210:1.8) -- cycle;
   \end{scope}
 
   % 季节标签
   \node[font=\footnotesize, green!60!black] at (90:1.4)  {春};
   \node[font=\footnotesize, red!60!black]   at (0:1.4)   {夏};
   \node[font=\footnotesize, orange!60!black] at (270:1.4) {秋};
   \node[font=\footnotesize, blue!60!black]  at (180:1.4) {冬};
 
   % 中心标签
   \node[font=\small, gray] at (0,0) {十二地支};
  % ??????? GitHub dark mode ??? SVG ???
  \begin{scope}[on background layer]
    \fill[white] (current bounding box.south west) rectangle (current bounding box.north east);
  \end{scope}
 \end{tikzpicture}
 \end{document}
